% ============================================================================
\section{GPAS Micro-Batch Allocation}
\label{sec:theory}
% ============================================================================

\subsection{Noise in optimizer coordinates}

AdamW divides coordinate $k$ of the gradient by $\sqrt{v_k}+\epsilon$, so
noise in coordinate $k$ affects the update in proportion to
$D_k=(\sqrt{v_k}+\epsilon)^{-1}$. We therefore measure gradient noise in these
coordinates, treating $D$ as fixed within a step and micro-batches as
independent. Let $e_i$ be the noise of one task-$i$ micro-batch and $V(m)$ the
noise of the accumulated gradient $A$:
\begin{equation}
    e_i=\mathbb E\bigl\|D(g_i-\mathbb E[g_i])\bigr\|_2^2,
    \qquad
    V(m)=\mathbb E\bigl\|D(A-\nabla F)\bigr\|_2^2
        =\sum_{i=1}^M\frac{w_i^2e_i}{m_i}.
    \label{eq:batch_variance}
\end{equation}
For a fixed student state and step size, the allocation enters the one-step
descent bound only through $V(m)$ (Appendix~\ref{app:descent}). Minimizing
$V(m)$ subject to $\sum_i m_i=G$ gives
\begin{equation}
    m_i\ \propto\ w_i\sqrt{e_i}.
    \label{eq:batch_gpas_allocation}
\end{equation}
We call this rule Gradient-Preconditioned Adaptive Sampling (GPAS). It is
Neyman allocation, the classical rule that samples each stratum in proportion
to its standard deviation, applied in the coordinates AdamW uses. The
continuous counts are clipped to $[m_{\min},m_{\max}]$, the unclipped ones
rescaled to sum to $G$, and the result rounded by largest remainder
(Appendix~\ref{app:allocation_proofs}).

The coordinates matter because raw and scaled noise can rank tasks
differently. In a controlled three-task example whose scaling reverses the
ranking, allocating by raw noise raises the scaled variance to $2.3\times$
that of Uniform, whereas GPAS lowers it to $0.68\times$
(Appendix~\ref{app:controlled}).

\subsection{Cost-GPAS}

Long-response tasks cost more per micro-batch, so the allocation with the
least variance per step need not have the least variance per hour. Let
$\tau_i$ be the measured time of one task-$i$ micro-batch and $C$ the fixed
time per step for synchronization, teacher loading, and the optimizer update,
so that a step takes $C+\sum_i m_i\tau_i$. A wall-clock budget $T$ then allows
$T/(C+\sum_i m_i\tau_i)$ steps, and the variance of the gradient averaged
over those steps is proportional to
\begin{equation}
    J_C(m)=\Bigl(C+\sum_{i=1}^M m_i\tau_i\Bigr)\,V(m).
    \label{eq:batch_cost_objective}
\end{equation}
Cost-GPAS minimizes $J_C$ over the feasible integer counts, a small
enumeration for four tasks. With $C=0$ the minimizer is
$m_i\propto w_i\sqrt{e_i/\tau_i}$ (Appendix~\ref{app:allocation_proofs}); as
$C$ grows relative to $\sum_i m_i\tau_i$, the minimizer approaches GPAS.

\subsection{Online estimates and implementation}

Every statistic for the next allocation comes from the step just completed.
Let $\bar g_i$ be the mean of task $i$'s micro-batch gradients in that step.
Then
\begin{equation}
    \widehat e_i=\frac{1}{m_i-1}\sum_{s=1}^{m_i}
    \bigl\|D(g_{i,s}-\bar g_i)\bigr\|_2^2
    \label{eq:noise_estimate}
\end{equation}
is the unbiased within-step estimate of $e_i$, and $m_{\min}=2$ makes it
available on every step. The allocation uses the current estimate without
smoothing. \missing{Confirm from the logged step-to-step change of
$\widehat e_i$; add smoothing only if the trace is unstable.} Response lengths
make timings noisy, so $\tau_i$ and $C$ are exponential averages of measured
values with decay $0.9$. We also log the predicted variance ratio relative to
Uniform,
\begin{equation}
    H=\frac{V(m^{\rm unif})}{V(m)},
    \label{eq:allocation_variance_ratio}
\end{equation}
which is at most $2$ because Uniform assigns $4$ micro-batches and
$m_i\leq8$. The first step uses Uniform because no statistics exist yet.

Computing $\widehat e_i$ needs one parameter-sized buffer for $\bar g_i$, a
few elementwise passes over the parameters per micro-batch, and one scalar
all-reduce; it adds no forward or backward pass. GPAS runs with conventional
AdamW. The allocation changes the noise term of AdamW's second moment, as any
change in batch composition does; Appendix~\ref{app:moment_target_proof}
bounds this change under the count limits and describes an optional
allocation-independent second-moment observation that the experiments do not
use. Algorithm~\ref{alg:batch_gpas} summarizes one step.

\begin{algorithm}[H]
\caption{One optimizer step with GPAS.}
\label{alg:batch_gpas}
\small
\begin{minipage}{0.97\linewidth}
\begin{enumerate}
    \item Choose bounded integer counts $m$ from the previous step's
    $\widehat e_i$ (and $\tau_i$, $C$ for Cost-GPAS); use Uniform on the first
    step.
    \item For every task $i$, take $m_ib$ new prompts, sample one response
    per prompt, obtain token-level teacher feedback, and accumulate each
    micro-batch gradient with weight $w_i/m_i$.
    \item From the same gradients and execution trace, compute
    $\widehat e_i$, update the timing averages, and log counts, losses,
    tokens, $H$, and timing.
    \item Apply one AdamW update.
\end{enumerate}
\end{minipage}
\end{algorithm}
